\chapter{Projection-Based Reduced Order Models}

Projection-based reduced order models (PROMs) can significantly reduce the time and storage required for the solution of high-fidelity (high-dimensional) discretized PDE-based computational models, also known as full-order models (FOMs). PROMs are particularly valuable in applications where time-critical models, such as digital twins, design optimization, and control problem optimization, are needed. These models project the high-dimensional system onto a lower-dimensional subspace, enabling rapid and efficient simulations while maintaining a high level of accuracy.

The creation of a PROM typically involves an offline-online decomposition strategy. In the offline phase, computationally intensive tasks are performed, such as generating a reduced-order basis (ROB) through methods like Proper Orthogonal Decomposition (POD) \cite{Sirovich1987, Balachandar1998}. In the online phase, the precomputed ROB is used to project the FOM onto the reduced subspace, allowing for efficient simulations \cite{Antoulas2005, Cuong05}. This approach has been shown to be effective for a variety of linear and nonlinear problems, particularly when the Kolmogorov n-width decays rapidly, as in the case of smooth solution manifolds \cite{Rowley2004, Benner2015}.

However, for problems exhibiting a large Kolmogorov n-width, such as those involving complex nonlinear phenomena, linear subspace approximations often struggle to capture the essential dynamics. This is particularly challenging in scenarios involving shock capturing, moving discontinuities, and other nonlinearities, where traditional linear techniques, like standard POD, fail to provide accurate representations. To overcome these limitations, nonlinear reduced-order modeling techniques have been developed, including local POD \cite{Grimberg2021,Bravo2024}, quadratic approximation models (QPROMs) \cite{Barnett2022}, neural network-augmented PROM (PROM-ANN) \cite{Barnett2023}, and nonlinear manifold approximations using dense and convolutional autoencoders \cite{Lee2020, Romor2023}.

% In this chapter, we aim to provide a comprehensive comparison of these state-of-the-art reduced-order modeling techniques, with a particular focus on moving discontinuity scenarios, such as those modeled by the Burgers' equation. We will explore various projection-based ROM methods, including Proper Orthogonal Decomposition (POD) \cite{Ares2023, Ares2024}, local-POD, Quadratic Approximation Manifolds, POD-ANN, and autoencoders. Furthermore, we will discuss both Galerkin and Least-Squares Petrov-Galerkin (LSPG) \cite{Carlberg2011, Carlberg2013, Carlberg2017} projection methods, highlighting their strengths and weaknesses in handling such complex dynamics.

% Additionally, we propose a novel deterministic approach based on Radial Basis Functions (RBFs), similar to the POD-ANN method, aimed at achieving comparable accuracy while maintaining computational efficiency and determinism. By providing a detailed analysis of these techniques in the context of complex scenarios, this paper serves as a guide for researchers and practitioners seeking the most suitable reduced-order modeling approach for their specific applications.
\section{Reduced Order Modeling (ROM)}

Let us consider a fully discrete set of governing equations in the form of a $\boldsymbol{\mu}-$parametric operator $\mathbf{R}:\mathbb{R}^{\mathcal{N}} \times \mathcal{P} \rightarrow \mathbb{R}^\mathcal{N}$, as

\begin{equation}
\mathbf{R}(\mathbf{U}; \boldsymbol{\mu}) = \boldsymbol{0} \ ,
\label{eq: residual discrete}
\end{equation}

\noindent where $\mathbf{U} \in \mathbb{R}^{\mathcal{N}}$ is the FOM solution vector containing the value of the solution on every degree of freedom of the spatial discretization, and $\boldsymbol{\mu} \in \mathcal{P} \subset \mathbb{R}^p$ is the parameter vector encapsulating, for example, geometrical variations, material properties, or boundary conditions.

For emphasis, let us reprint the iteration method used to obtain the solution at the next step, given the FOM solution at step $n$, symbolically represented as $\mathbf{U}_n$:

\begin{subequations}
\begin{align}
\mathbf{J}(\mathbf{U}_{n+1}^k) \delta \mathbf{U} & = \mathbf{R} (\mathbf{U}_{n+1}^{k}; \boldsymbol{\mu})\\
\mathbf{U}_{n+1}^{k+1} &= \mathbf{U}_{n+1}^{k} + \delta\mathbf{U},
\end{align}
\end{subequations}
\noindent where:
\begin{itemize}
\item $\mathbf{U} \in \mathbb{R}^{\mathcal{N}}$ is the solution vector at the current step.
\item $\mathbf{J} = \frac{\partial \mathbf{R}}{\partial \mathbf{U}} \in \mathbb{R}^{\mathcal{N} \times \mathcal{N}}$ is the Jacobian matrix (or its approximation).
\item $\delta \mathbf{U} \in \mathbb{R}^{\mathcal{N}}$ is the differential increment.
\item $k$ is the current iterate.
\end{itemize}

Solving this linear system of equations can be computationally expensive, especially for large-scale problems where the dimension \(\mathcal{N}\) is very high. This is where reduced-order modeling becomes crucial, as it aims to achieve significant computational efficiency by reducing the dimensionality of the problem while preserving the essential dynamics of the original system.

\subsection{Manifold Projection-Based ROMs}
\label{subsec: manifold-projection-based-rom}

We begin by approximating the solution state variable through a general decoder function defined as:
\begin{equation}
\mathbf{U} \approx D_u(\mathbf{q}),
\label{eq: Genral Decoder}
\end{equation}
where $D_u: \mathbb{R}^r \rightarrow \mathbb{R}^{\mathcal{N}}$ maps $\mathbf{q}$ to $\mathbf{U}$.

Correspondingly, we introduce its encoder:
\begin{equation}
\begin{aligned}
E_u: \ & \mathbb{R}^{\mathcal{N}} \rightarrow \mathbb{R}^r \
& \mathbf{U} \mapsto \mathbf{q}.\
\end{aligned}
\end{equation}
When Eq.\ref{eq: Genral Decoder} is substituted into Eq.\ref{eq: residual discrete}, we obtain a residual characterized by:
\begin{equation}
\mathbf{R}(\mathbf{U};\boldsymbol{\mu})=\mathbf{R}(D_u(\mathbf{q});\boldsymbol{\mu}).
    \label{eq: Discrete residual (decoder)}
\end{equation}
We now aim to minimize the non-linear residual in some norm $\mathbf{G}$ in order to meet the minimum-residual optimality of the projection approximation. Specifically, we seek to solve the following minimization problem:
\begin{equation}
    \min_{\mathbf{q} \in \mathbb{R}^n} \frac{1}{2}\left\lVert \mathbf{R}(D_u(\mathbf{q});\boldsymbol{\mu}) \right\rVert_{\mathbf{G}}^2 
    \label{Minimization of the residual}
\end{equation}
where $\|\cdot\|_{\mathbf{G}}$ denotes a norm defined by a symmetric, positive definite (SPD) matrix $\mathbf{G} \in \mathbb{R}^{N\times N}$. Alternatively, we can write the above problem as follows:
\begin{equation}
    \min_{\mathbf{q} \in \mathbb{R}^n} \frac{1}{2}\mathbf{R}(D_u(\mathbf{q}))^T\mathbf{G}\mathbf{R}(D_u(\mathbf{q})).
\end{equation}
By satisfying the minimum-residual optimality of Eq.\ref{Minimization of the residual}, a general relationship between projection-based reduced-order models and minimum-residual reduced-order models can be established, as follows: 
\begin{equation}
\begin{aligned}
    \left[\frac{\partial \mathbf{R}}{\partial \mathbf{U}} \frac{\partial (D_u(\mathbf{q}))}{\partial \mathbf{q}}\right]^T\mathbf{G}\mathbf{R}&=\mathbf{0}\\
    \left[\mathbf{J}\frac{\partial (D_u(\mathbf{q}))}{\partial \mathbf{q}}\right]^T\mathbf{G}\mathbf{R}&=\mathbf{0}.
\end{aligned}
\end{equation}

This version of the projected residual can be seen as a general encoder of the residual, that is:
% \begin{equation}
%     E_r (\mathbf{R}) = \mathbf{\Psi}^T \mathbf{R} = \frac{\partial (D_u(\mathbf{q}))}{\partial \mathbf{q}}^T \mathbf{J}^T\mathbf{G} \mathbf{R}
%     \label{Encoder of the residual}
% \end{equation}
\begin{equation}
    E_r (\mathbf{R}) = \frac{\partial (D_u(\mathbf{q}))}{\partial \mathbf{q}}^T \mathbf{J}^T\mathbf{G} \mathbf{R}
    \label{Encoder of the residual}
\end{equation}

We can then propose to minimize the following expression depending on small dimensional quantities $\mathbf{r} \in \mathbb{R}^\ell$, and $\mathbf{q} \in \mathbb{R}^k$:

\begin{equation}
    \mathbf{r}(\mathbf{q}; \boldsymbol{\mu})= E_r\left(\mathbf{R}(D_u(\mathbf{q}) ; \boldsymbol{\mu}) \right) = \boldsymbol{0}
\end{equation}

\noindent where $\mathbf{r}(\mathbf{q}; \boldsymbol{\mu}) : = E_r\left(\mathbf{R}(D_u(\mathbf{q}\right) ; \boldsymbol{\mu}) )$. Moreover, given a ROM solution at a step $n$, symbolically represented as $\tilde{\mathbf{U}}_n \in \mathbb{R}^{\mathcal{N}}$, the solution corresponding to the next step can be obtained via a iteration method like

\begin{equation}
    \frac{\partial{\mathbf{r}}}{\partial{\mathbf{q}}} \delta \mathbf{q}=-\mathbf{r},
    \label{iteration method}
\end{equation}

\begin{subequations}
\begin{align}
\frac{\partial E_r(\mathbf{R})}{\partial \mathbf{R}} 
\frac{\partial \mathbf{R}}{\partial \mathbf{U}}\frac{\partial D_u(\mathbf{q})}{\partial \mathbf{q}} \delta \mathbf{q}^{k} & =   - E_r\left(\mathbf{R} \right)
\label{eq: Newton Raphson Enoder Residual}  \\
\mathbf{q}^{k+1} & = \mathbf{q}^{k} + \delta\mathbf{q} \ ,
\end{align}
\end{subequations}

Elaborating on equation \ref{eq: Newton Raphson Enoder Residual} given the definition of the encoder of the residual in eq. \ref{Encoder of the residual}, we obtain:
\begin{equation}
    \frac{\partial (D_u(\mathbf{q}))}{\partial \mathbf{q}}^T \mathbf{J}^T\mathbf{G} \mathbf{J} \frac{\partial D_u(\mathbf{q})}{\partial \mathbf{q}} \delta \mathbf{q}^{k} =   - \frac{\partial (D_u(\mathbf{q}))}{\partial \mathbf{q}}^T \mathbf{J}^T\mathbf{G} \mathbf{R}
    \label{eq:prom_general}
\end{equation}

\subsubsection{Galerkin Projection}
The Galerkin projection is particularly effective when dealing with symmetric positive definite (SPD) Jacobians. In such scenarios, it yields search directions that are minimum-residual optimal. This optimality is achieved with $\mathbf{G}=\mathbf{J}^{-T}$, leading to the following formulation:

\begin{equation}
\frac{\partial (D_u(\mathbf{q}))}{\partial \mathbf{q}}^T \mathbf{J} \frac{\partial D_u(\mathbf{q})}{\partial \mathbf{q}} \delta \mathbf{q}^{k} = - \frac{\partial (D_u(\mathbf{q}))}{\partial \mathbf{q}}^T \mathbf{R}
\label{eq:prom_general_galerkin}
\end{equation}

\subsubsection{Least-Squares Petrov-Galerkin (LSPG) Projection}
In nonlinear problems, Jacobians often are not SPD. The LSPG method addresses this by setting $\mathbf{G} = \mathbf{I}$, the identity matrix. This choice ensures that the LSPG method maintains the minimum-residual property, crucial for the stability and accuracy of the solution. The projection equation in this context is:

\begin{equation}
\frac{\partial (D_u(\mathbf{q}))}{\partial \mathbf{q}}^T \mathbf{J}^T \mathbf{J} \frac{\partial D_u(\mathbf{q})}{\partial \mathbf{q}} \delta \mathbf{q}^{k} = - \frac{\partial (D_u(\mathbf{q}))}{\partial \mathbf{q}}^T \mathbf{J}^T \mathbf{R}
\label{eq:prom_general_lspg}
\end{equation}


\section{Linear PROMs}
\section{Piecewise linear PROMs}
\section{Quadratic PROMs}
\section{Nonlinear manifold PROMs}
\section{Nonlinear map augmented PROMs}